% Einheit 4 von 5 · Binomische Formeln (binomische-formeln.md, Einheit 2) – Hauptnummern 28–38
\setcounter{aufgabe}{27}
\zweigkopf{Einheit 4 von 5 · Binomische Formeln}{Hier lernst du, Quadrate von Klammern mit den drei binomischen Formeln auszumultiplizieren · neu in diesem Jahr · P10 · baut auf: Summe mal Summe (Einheit 3), Minusklammer (Einheit 1), Quadratzahlen}{e4}

\begin{aufgabe}{Ich erkenne, ob eine Klammer mit sich selbst malgenommen wird. Kreuze an, ob der Term eine Klammer ist, die mit sich selbst malgenommen wird. Wenn ja, schreibe ihn in der anderen Form auf: als Klammer mal Klammer oder als Quadrat einer Klammer. Multipliziere nichts aus.}
\begin{geruest}
\gz{a}{$(x + 8)^2$}{\janein \quad $=$ \leerfeld}
\gz{b}{$(x + 4) \cdot (x - 4)$}{\janein \quad $=$ \leerfeld}
\gz{c}{$(m - 6)^2$}{\janein \quad $=$ \leerfeld}
\gz{d}{$(y + 2) \cdot (2 + y)$}{\janein \quad $=$ \leerfeld}
\gz{e}{$(2x + 1) \cdot (x + 1)$}{\janein \quad $=$ \leerfeld}
\end{geruest}
\end{aufgabe}

\begin{aufgabe}{Ich erkenne, was $a$ und was $b$ ist. In $(a + b)^2$ und $(a - b)^2$ ist $a$ das erste Glied in der Klammer und $b$ das zweite, ohne sein Vorzeichen. Schreibe auf, was $a$ und was $b$ ist. Rechne nichts aus.}
\begin{geruest}
\gz{a}{$(x + 9)^2$}{\feld{a}\feld{b}}
\gz{b}{$(y - 4)^2$}{\feld{a}\feld{b}}
\gz{c}{$(4x + 1)^2$}{\feld{a}\feld{b}}
\gz{d}{$(3m - 2n)^2$}{\feld{a}\feld{b}}
\gz{e}{$(6 + z)^2$}{\feld{a}\feld{b}}
\end{geruest}
\end{aufgabe}

\begin{aufgabe}{Ich erkenne, welche binomische Formel passt. Plus in der Klammer und Quadrat: erste Formel. Minus in der Klammer und Quadrat: zweite Formel. Plus-Klammer mal Minus-Klammer mit denselben Gliedern: dritte Formel. Sonst: keine. Schreibe auf, welche Formel passt. Rechne nichts aus.}
\begin{geruest}
\gz{a}{$(x + 12)^2$}{\leerfeld}
\gz{b}{$(m + 7) \cdot (m - 7)$}{\leerfeld}
\gz{c}{$(y - 5) \cdot (y + 6)$}{\leerfeld}
\gz{d}{$(z - 10)^2$}{\leerfeld}
\gz{e}{$(u + 2) \cdot (u + 3)$}{\leerfeld}
\end{geruest}
\end{aufgabe}

\begin{aufgabe}{Ich kann eine Klammer mit einer binomischen Formel ausmultiplizieren. Multipliziere mit der passenden binomischen Formel aus.}
\beispiel{(x + 6)^2 &= x^2 + 2 \cdot x \cdot 6 + 6^2 \\ &= x^2 + 12x + 36}
\begin{gleichungsraster}[2]
\gl{(x + 1)^2} & \gl{(x + 4)^2} \\ \noalign{\vspace{5mm}}
\gl{(m + 2)^2} & \gl{(y + 10)^2} \\ \noalign{\vspace{5mm}}
\gl{(x - 9)^2} & \gl{(x + 8) \cdot (x - 8)} \\ \noalign{\vspace{5mm}}
\end{gleichungsraster}
\end{aufgabe}

\begin{aufgabe}{Ich kann eine Klammer mit einer binomischen Formel ausmultiplizieren – weiter. Entscheide, ob eine binomische Formel passt, und multipliziere aus.}
\begin{gleichungsraster}[2]
\gl{(4 - z) \cdot (4 + z)} & \gl{(x + 6) \cdot (x - 1)} \\ \noalign{\vspace{5mm}}
\end{gleichungsraster}
Multipliziere aus. Die Vorzahl gehört zu $a$ oder $b$ und wird mit quadriert.
\begin{gleichungsraster}[2]
\gl{(3x + 2)^2} & \gl{(2m - n)^2} \\ \noalign{\vspace{5mm}}
\end{gleichungsraster}
\end{aufgabe}

\begin{aufgabe}{Ich kann eine Klammer mit einer binomischen Formel ausmultiplizieren – weiter. Multipliziere erst die Klammer mit der Formel aus, dann den Rest. Fasse zusammen.}
\begin{gleichungsraster}[2]
\gl{4 \cdot (x - 1)^2} & \gl{-(x - 7)^2} \\ \noalign{\vspace{5mm}}
\gl{(x + 2)^2 - 4x} & \\ \noalign{\vspace{5mm}}
\end{gleichungsraster}
Zeige durch Umformen, dass der Term $(x - 6)^2 - 11$ gleich $x^2 - 12x + 25$ ist. (P10 2017)
\begin{gleichungsraster}[3]
\gl{(x - 6)^2 - 11} & \\ \noalign{\vspace{5mm}}
\end{gleichungsraster}
\end{aufgabe}

\begin{aufgabe}{Ich kann mit einer binomischen Formel im Kopf rechnen. Schreibe die Zahl als Summe oder Differenz mit einer glatten Zahl und rechne mit der passenden Formel.}
\beispiel{41^2 &= (40 + 1)^2 = 1600 + 80 + 1 = 1681}
\begin{teilezwei}
\tz $31^2 =$ \leerfeld & \tz $19^2 =$ \leerfeld \\
\tz $52 \cdot 48 =$ \leerfeld & \tz $99^2 =$ \leerfeld \\
\end{teilezwei}
\end{aufgabe}

\begin{aufgabe}{Ich finde den Fehler bei einer binomischen Formel. Finde den Fehler, benenne ihn und korrigiere die Rechnung. Rechne danach die Aufgabe darunter selbst.}
\begin{teile}
\teil Tim rechnet so:
\end{teile}
\rechnung{(x + 11)^2 &= x^2 + 121}
\begin{gleichungsraster}[2]
\gl{(m + 13)^2} & \\ \noalign{\vspace{5mm}}
\end{gleichungsraster}
\begin{teile}
\teil Sara rechnet so:
\end{teile}
\rechnung{(x - 10)^2 &= x^2 + 20x + 100}
\begin{gleichungsraster}[2]
\gl{(y - 8)^2} & \\ \noalign{\vspace{5mm}}
\end{gleichungsraster}
\end{aufgabe}

\begin{aufgabe}{Ich finde den Fehler bei einer binomischen Formel – weiter. Finde den Fehler, benenne ihn und korrigiere die Rechnung. Rechne danach die Aufgabe darunter selbst.}
\begin{teile}
\teil Max rechnet so:
\end{teile}
\rechnung{(5x + 2)^2 &= 5x^2 + 20x + 4}
\begin{gleichungsraster}[2]
\gl{(4y + 3)^2} & \\ \noalign{\vspace{5mm}}
\end{gleichungsraster}
\begin{teile}
\teil Jana rechnet so:
\end{teile}
\rechnung{-(x + 6)^2 &= -x^2 + 12x + 36}
\begin{gleichungsraster}[2]
\gl{20 - (x + 4)^2} & \\ \noalign{\vspace{5mm}}
\end{gleichungsraster}
\end{aufgabe}

\begin{aufgabe}{Ich kann begründen, warum $(a + b)^2$ nicht $a^2 + b^2$ ist.}
\begin{teile}
\teil Rechne $(6 + 4)^2$ und $6^2 + 4^2$. Was fällt dir auf?
\teil Das Quadrat im Bild hat die Seitenlänge $x + 2$. Schreibe in jedes der vier Teile seinen Flächeninhalt. Begründe damit, dass $(x + 2)^2 = x^2 + 4x + 4$ ist.
\end{teile}

\flaechenbild{3.5}{1.5}{3.5}{1.5}{x}{2}{x}{2}

\begin{teile}
\teil Gibt es Zahlen $a$ und $b$, für die $(a + b)^2 = a^2 + b^2$ stimmt? Begründe.
\end{teile}
\end{aufgabe}

\begin{aufgabe}{Ich kann mit einer binomischen Formel die Fläche eines Rahmens berechnen. Ein quadratisches Foto mit der Seitenlänge $q$~cm bekommt ringsum einen $2$~cm breiten Rahmen.}
\begin{teile}
\teil Gib einen Term für die Fläche von Foto und Rahmen zusammen an und multipliziere aus.
\teil Gib einen Term für die Fläche des Rahmens allein an.
\teil Das Foto ist $10$~cm breit. Wie groß ist die Fläche des Rahmens?
\end{teile}
\end{aufgabe}
